一阶语言的\emph{语义}由该语言的一个\emph{结构}给出。
它包含一个\emph{论域}，并为每个常项符号指派该论域中的元素。
函数符号由函数来解释，关系符号则由论域上的关系来解释。
从变元集合到论域的一个函数称为一个\emph{变元指派}。
\emph{满足关系}将结构、变元指派和公式联系起来；
$\Sat{M}{!A}[s]$ 通过对~$!A$ 的结构进行归纳来定义。
$\Sat{M}{!A}[s]$ 仅依赖于实际出现在~$!A$ 中的符号的解释，
特别地，如果 $!A$ 不含自由变元，则不依赖于 $s$。
因此，如果 $!A$ 是一个语句，那么只要 $\Sat{M}{!A}[s]$ 对任意（或所有）~$s$ 成立，就有 $\Sat{M}{!A}$。

满足关系是所有语义概念的基础。
一个语句是\emph{有效的}，记作 $\Sat{{}}{!A}$，如果它在每个结构中都被满足。
语句集~$\Gamma$ \emph{语义地后承}语句 $!A$，记作 $\Gamma \Entails !A$，
当且仅当所有满足~$\Gamma$ 中每一语句的 $\Struct{M}$ 都满足 $!A$。
集合~$\Gamma$ 是\emph{可满足的}，当且仅当存在一个结构满足~$\Gamma$ 中的每一语句，否则就是\emph{不可满足的}。
这些概念相互关联，例如，$\Gamma \Entails !A$ 当且仅当 $\Gamma \cup \{\lnot !A\}$ 不可满足。